\documentclass[12pt, letterpaper]{article}

\usepackage[utf8]{inputenc}
\usepackage[T1]{fontenc}

\usepackage{amsmath, amssymb, amsfonts}

\usepackage{graphicx}
\usepackage{xcolor}


\usepackage[english]{babel}   
\usepackage{float}   
\usepackage{url}    
\usepackage{hyperref}

\hypersetup{colorlinks=true,linkcolor=blue,urlcolor=blue}
\urlstyle{rm}

\usepackage{geometry}
\geometry{margin=1in}

\title{DAT158 End-to-end ML project}
\author{Thobias Kvalvik H{\o}ivik}
\date{}


\begin{document}

\maketitle

\section{The idea}

My idea for the project was to make 
a (very simplified) physical simulation where 
you have a spacecraft, unaffected by gravity, 
whose goal is to move somewhere. Crucially, the 
only things the spacecraft should be allowed to do 
are rotate, and choose how much thrust is applied. 

This idea was inspired by 
\href{https://www.youtube.com/watch?v=L8XkPfY-jc8}{this old video by Pezza's Work}
which I have always wanted to try to recreate. 

Now, the most natural approach, which is 
the same used by said creator, is to use a small neural network, 
which to my understanding lies a bit outside the scope 
of the syllabus, but as the solution I landed 
on is a form of reinforcement learning I figured 
it'd be fine. 

The demo may be found on \href{https://thobiashoivik.com/lab/rocket-reinforcement-learning}{my personal website}.

I won't devote an entire section to 
the simulation, but it is quite simple. 
The controller chooses thrust and torque and the 
simulation updates how those change the 
motion of the rocket. 
The rocket has a position, velocity, rotation, 
and angular velocity. 
They are updated each $\Delta t = 1/30$ seconds.
Thrust adds acceleration of up to $4$ m/s$^2$ along the nose,
and torque adds angular acceleration of up to $6$ rad/s$^2$.
Each step also multiplies linear velocity by $0.99$ and
angular velocity by $0.94$. This artificial damping makes
both motion and spinning gradually slow down.

\section{How does it work?}

In supervised learning, we might have 
an input $x$, a known answer $y$, and a prediction 
$\widehat y$ with a loss such as 
$(\widehat y - y)^2$ which tells you how to 
adjust the model. 

In our problem, where we are 
simulating a rocket in space, 
we want the model to have access 
to the current rocket's situation with 
the model then choosing thrust and torque. 

We do not have a predetermined answer to what the 
correct answer should be. We may only 
run the simulation and judge what happens next. 
This also creates a slight complication, that 
being that an action changes what 
situation the model will encounter next. 

The \textbf{policy} is the model's rule for 
choosing actions from observations. 
This implementation uses two networks, 
\begin{itemize}
    \item An actor, who chooses how to control the rocket, and 
    \item A critic, who estimates how much future reward 
        the current situation is likely to produce. 
\end{itemize}

The critic's job is to help train the actor, and 
does not issue any steering commands. 

In each simulation step, 
the actor is given seven numbers 
\[
    s_t = 
    \begin{bmatrix}
        (g_x - x)/ 3 \\ 
        (g_y - y)/ 3 \\ 
        v_x/3 \\ 
        v_y/3 \\ 
        \sin \theta \\ 
        \cos \theta \\ 
        \omega/6 
    \end{bmatrix}
\]
where the first two entries correspond to 
the signed horizontal and vertical offsets from the goal, the next two 
represent velocity at this moment, those after 
that correspond to rotation and the last 
is the angular velocity (how much the rocket is spinning). 

Positions and velocities are scaled 
to keep their ranges manageable and heading/rotation 
uses sine and cosine since we 
don't want to deal with $0 \equiv 2\pi$ wrap-around 
shenanigans. 
All entries except sine and cosine are also clipped to $[-2,2]$.
The network does not receive the target's velocity.

The actor is then a small network with $7$ inputs, 
two hidden layers with $32$ neurons each, and 
$2$ outputs. This is pretty much identical to the 
architecture used in the aforementioned YouTube 
video, except that there the two hidden layers 
have only $9$ neurons.

Each hidden neuron in our network computes a 
weighted sum, adds a bias and applies $\tanh$ which 
lets the network represent nonlinear relationships so
the appropriate thrust can depend jointly 
on a bunch of observations in a 
potentially weird way. 

It is worth mentioning that 
during training, the actor describes 
a distribution of possible actions and its 
two outputs are means, written $\mu$, for 
two Gaussian distributions. We sample them such: 
\[
    u = \mu + \sigma \epsilon, \ \epsilon \sim \mathcal N(0,1)
\]
where 
\begin{itemize}
    \item $\mu$ is the network's preferred output, 
    \item $\sigma$ controls the amount of variation, and 
    \item $\epsilon$ supplies some new random variation.
\end{itemize}

The implementation also 
learns the two standard deviations, which are shared 
across observations, whereas 
the means change depending on what is required 
given the rocket's current state. 

Sampling provides a sort of 
exploration. If the model always repeated 
its current preferred action it would 
collect little evidence about alternatives. 
A stochastic policy tries "nearby" alternatives 
and can potentially discover some 
which work better. 

This approach is described in 
\href{https://spinningup.openai.com/en/latest/spinningup/rl_intro.html}{Spinning Up's Introduction}.

The sampled numbers $u$ are 
latent outputs and they can, 
annoyingly, be arbitrarily large so we convert 
them into bounded controls, 
\[
    T = \frac{\tanh(u_0) + 1}{2}, \ \tau = \tanh(u_1).
\]
This gives thrust $T$ between $0$ and $1$ and 
torque $\tau$ between $-1$ and $1$. 

After applying these controls, the 
simulation produces a reward. 
The scoring rule rewards 
approaching and remaining near the target, and 
penalizes distance, excessive speed near the target, 
"fuel use", torque and spinning. 

More precisely, writing $d_t$ and $d_{t+1}$ for the distances
before and after a physics step with the same target held fixed,
the reward is
\begin{align*}
    q_t &= \exp(-2d_{t+1}^2),\\
    r_t &= 4(d_t-d_{t+1}) - 0.10d_{t+1} + 0.15q_t\\
        &\quad - 0.08q_t\|\mathbf v_{t+1}\|^2 - 0.008T_t\\
        &\quad - 0.004\tau_t^2 - 0.002\omega_{t+1}^2.
\end{align*}
An escape adds a penalty of $-1$. "Fuel use" is an engine-use
penalty, rather than an actual tank that empties.

This works reasonably well, but we will discuss 
some annoying side effects from this metric later on. 

Immediate reward alone is not sufficient. 
If the rocket approaches the target quickly, 
braking might briefly increase its distance 
from the target, but prevent a large overshoot. An 
algorithm that judged only the next 
instance could discourage useful braking. 

We therefore care about return, i.e. accumulated 
future reward. The implementation multiplies the contribution
by $0.99$ for each additional step, 
\[
    G_t = r_t + 0.99r_{t+1} + 0.99^2r_{t+2} + \cdots
\]
Here $r_t$ is the immediate reward and $G_t$ 
represents longer-term outcome. 

The critic's job, as mentioned, is to 
predict the longer-term outcome. 
We denote the prediction $V(s)$, where $s$ is 
the observation. For instance, 
it might learn that "close to the target, but 
quickly moving away" is less promising 
than "slightly farther away but moving slowly toward 
it". 

This leads to an advantage, i.e. an estimate 
of how favorable the sampled decision was 
relative to what the critic expected. 

The first term is a temporal-difference error, 
\[
    \delta_t = r_t + 0.99 V(s_{t+1}) - V(s_t),
\]
which we can read as the reward 
received, in addition to the 
estimated value of the next state, minus 
the value we previously expected. 

The implementation combines error from 
successive steps using Generalized Advantage 
Estimation, 
\[
    A_t = \delta_t + (0.99 \cdot 0.95)A_{t+1}. 
\]

Repeatedly expanding this expression gives 
the current error plus progressively smaller 
contributions from later errors. It 
lets later consequences influence the assessment 
of an earlier action. 
The $0.95$ parameter balances reliance on short-term 
critic predictions against reliance on 
longer, noisier sequences of experience. 
That is the purpose of 
\href{https://arxiv.org/pdf/1506.02438}{GAE}. 

These equations describe an uninterrupted trajectory.
An escape removes the next-state value term. A time limit
retains that term using the final observation, but the advantage
recursion stops across either kind of episode reset.

These are raw advantages. Before updating the actor, they are
centered using the whole batch's mean and divided by its standard
deviation (with a small numerical safeguard). I will write the
normalized advantages as $\widetilde A_t$. Centering can change
an individual sign: a positive normalized advantage means
favorable relative to the batch average. The critic still uses
the raw advantages.

Now we can elucidate how the critic 
helps guide the actor. 

We save the action the actor actually sampled.  
During optimization, we ask how likely 
the current actor considers the same action, 
\begin{itemize}
    \item a favourable normalized advantage encourages increasing its likelihood, while 
    \item an unfavourable normalized advantage encourages decreasing its likelihood.
\end{itemize}

Since the outputs are continuous (as a theoretical compsci person 
I would say not really, but it's a fine assumption), 
likelihood means probability density. 
We make actions near the sampled 
value more or less likely. 

The basic policy-gradient learning signal involves 
\[
    \widetilde A_t \log \pi_w(u_t \mid s_t), 
\]
where $\pi_w$ is the action distribution produced 
by actor weights $w$. 
The density is already differentiable; its logarithm gives a useful
gradient for changing the relative likelihood of sampled outputs.
Backpropagation computes how changing each weight 
affects this expression, and the optimizer 
can update the weights accordingly. 

We do not backpropagate through the rocket's 
physics; instead the simulator 
supplies experience, while 
the gradient is computed through the actor's 
likelihood calculation. 
The \href{https://spinningup.openai.com/en/latest/spinningup/rl_intro3.html}{policy-gradient derivation}
explains why this can estimate how to improve 
expected return. 

For one Gaussian output, the derivative of log-likelihood 
with respect to its mean is 
\[
    \frac{\partial \log \pi(u)}{\partial \mu} = 
    \frac{u-\mu}{\sigma^2}. 
\]

If the sampled output was above the mean, 
a positive normalized advantage tends to pull 
the mean upward towards it, while 
a negative normalized advantage tends to push the mean further 
away. The full network update combines many 
such samples, so an individual sample does not 
determine the result by itself. 

The algorithm used, Proximal Policy Optimization, 
or PPO, adds a restraint. We reuse a batch 
for several passes, but that batch 
comes from the actor before the updates. 
Large changes would make the collected 
experience less representative of the current policy. 

PPO compares the new and old densities for each recorded 
action, 
\[
    \rho_t = \frac{\pi_{\text{new}}(u_t\mid s_t)}{\pi_{\text{old}}(u_t\mid s_t)}. 
\]

A ratio of $1.1$ means its density increased by $10\%$. 
With the clipping setting used, 
\begin{itemize}
    \item for positive normalized advantage, this sample 
    stops providing further encouragement once its ratio exceeds 
    $1.2$, while  
    \item for negative normalized advantage, it stops providing further 
    encouragement to decrease its 
    density once the ratio falls below $0.8$.
\end{itemize}
This removes the incentive for further 
favourable movement from that 
sample's part of the objective. 
Other samples can still change the same weight. 
This behavior is explained in 
\href{https://spinningup.openai.com/en/latest/algorithms/ppo.html}{the PPO documentation}. 

We then fit the critic with squared error against the targets
\[
    Y_t = V_{\mathrm{old}}(s_t) + A_t,
\]
using the raw advantages. These targets are generated from rewards
and existing value predictions, and are kept fixed during the five
optimization passes.

The training loop in the training 
script is 
\begin{enumerate}
    \item Run $128$ rocket simulations, collecting 
    $128$ steps for each, giving $16,384$ transitions. 
    \item Save observations, sampled outputs, rewards, 
    old action densities, and critic predictions. 
    \item Calculate raw advantages 
    and critic targets, then normalize the 
    advantages. 
    \item Shuffle the batch and make five optimization 
    passes, using groups of $1,024$ transitions. 
    \item Update the actor with PPO and the critic 
    with regression. 
    \item Collect new experience with the updated 
    actor and repeat. 
\end{enumerate}

\section{Actual implementation and such}

The network's derivatives and 
optimizer were implemented directly with 
NumPy.
The optimizer, Adam, uses recent gradient history to 
adjust parameter updates and a small 
entropy bonus encourages the latent action 
distributions to retain some variation. 

\paragraph{Training experience and settings}
There is no downloaded dataset of correct steering commands.
The simulation generates observations, sampled actions, rewards
and following observations. Episodes last $240$ steps, or $8$
seconds, unless the rocket escapes beyond $|x|>5$ or $|y|>5$.
Starting positions, velocities, headings and targets are randomized.
Targets are drawn from a mixture of static positions and smooth
sinusoidal paths. During the first $20\%$ of training, the reset
conditions and target motion gradually become more difficult.
The handwritten controller supplies no training examples.

The final run collected $30,015,488$ transitions with seed $17$.
The initial actor learning rate was $0.0004$, linearly decreasing
to about $30\%$ of that rate; the critic used twice the actor rate.
PPO clipping was $0.2$, and the latent Gaussian entropy coefficient
was $0.003$. These settings and the earlier batch sizes are recorded
alongside the weights in \path{public/models/rocket-rl.json}.

Finally, the browser runs a saved actor. At 
each step it reads the seven inputs, evaluates 
the network, and converts its means into thrust and 
torque. It turns off the training-time random 
sampling. It does not learn anything when 
interacted with in the page, and it does 
not search through possible future trajectories. 

The network was trained in about $4$ minutes on my 
own desktop computer with specs:
\begin{itemize}
    \item AMD Ryzen $7$ $7800$X$3$D, 
    \item $64$gb of $5600$MHz DDR$5$, 
    \item NVIDIA GeForce RTX $4080$ SUPER (not really relevant).
\end{itemize}
The measured runtime was $234.27$ seconds on the CPU;
the GPU was not used.

There is also a handwritten controller explicitly 
containing rules for pulling toward the target, 
damping velocity, and turning toward the 
desired acceleration. 

\paragraph{Evaluation}
Each controller was evaluated on the same $100$ seeded episodes
of $240$ steps, using seed $90210$ and randomized static or
sinusoidal targets. These scenarios were inspected during
development, so the results are validation results rather than
an untouched final test.

Success means a mean distance below $0.45$ m and mean speed below
$0.8$ m/s during the final $60$ steps. After an escape, remaining
steps count as $\sqrt{50}$ m of tracking error, a reward of
$-0.1\sqrt{50}$ per step, and zero engine use and spin.
This keeps early failure from artificially improving the distance
metric; low average thrust alone does not mean successful flight.

\begin{table}[htbp]
\centering
\small
\setlength{\tabcolsep}{5pt}
\begin{tabular}{lrrrrr}
\hline
Controller & Distance (m) & Thrust & Spin (rad/s) & Return & Success\\
\hline
Random      & 5.548 & 0.166 & 0.087 & -151.327 & $0\%$\\
Untrained   & 5.566 & 0.165 & 0.017 & -151.674 & $0\%$\\
Early PPO   & 0.609 & 0.141 & 2.669 & 9.958    & $96\%$\\
Middle PPO  & 0.397 & 0.196 & 2.989 & 17.947   & $100\%$\\
Final PPO   & 0.400 & 0.167 & 2.349 & 19.715   & $100\%$\\
Handwritten & 0.433 & 0.075 & 1.215 & 21.128   & $99\%$\\
\hline
\end{tabular}
\caption{Mean validation results. Spin is absolute angular speed;
thrust is on a scale from $0$ to $1$.}
\end{table}

Checkpoints were saved at $0$, $3,014,656$, $12,009,472$ and
$30,015,488$ transitions. The learning curve shown on the lab page
uses separate $32$-episode evaluations during training.

\paragraph{Reproducibility}
The root \texttt{README.md} describes how to install dependencies,
train and evaluate the controller, and build and run the website.
The training scripts and their detailed documentation are under
\texttt{scripts/rocket-rl}. Training requires Python $3.11$ or newer
and NumPy $2.3.5$. From the repository root:
\begin{verbatim}
python -m pip install -r scripts/rocket-rl/requirements.txt
python scripts/rocket-rl/train.py --steps 30000000 --seed 17
python scripts/rocket-rl/evaluate.py --verify
\end{verbatim}
The website requires Node.js $22.12$ or newer. Its build and
deployment instructions, saved model and simulation source are
included in the repository. The implementation was also checked
with finite-difference gradients and browser/Python parity tests.

Since the project was developed first and 
foremost with 
the goal of putting it on my website, 
the mangled remains of the project I have put here 
may be impractical to work with. 
Sorry for this.  

\section{Discussion}

The end result is a pretty neat toy 
to play around with and it performs pretty well, 
being able to reach the target quite effectively, 
supposing you are using the final policy. 

Perhaps unsurprisingly, the hand-written approach looks
a lot smoother. The learned controller has slightly lower
mean tracking distance, but the handwritten controller uses
less thrust, spins less and actually receives higher mean reward.
So the learned policy is not simply better according to the
reward metric. This points to our maybe not having fine-tuned the
reward metric well enough for the kind of flight I wanted.
It could also reflect PPO not fully optimizing the reward
that was already defined.

A major issue with the model is that if you very 
rapidly move the target around it is very easy for 
the rocket to end up flying too far off the screen 
as the model wants to correct by doing a really 
wide turn, which isn't great.
In fact, putting the goal in the bottom right or 
top right 
makes the model send the rocket off the screen every time 
with the latest policy.  

These are observations from interacting with the cursor, rather
than results from the fixed validation scenarios above. The target
positions and sudden movements can differ from the training
distribution, and the seven inputs do not include target velocity.
Also, the reset boundary lies beyond the visible flight area, so
leaving the screen is not the same as reaching an episode escape.
The single training seed and these failures leave plenty of room
for improvement. More varied target paths, stronger penalties for
unwanted spinning and evaluation across additional seeds would
be sensible next experiments; they have not been carried out here.



\paragraph{AI use declaration}


OpenAI's ChatGPT was, in part, used to converge 
on certain design choices, along with technical 
explanations of how such choices are implemented, 
as well as parts of actual implementation, 
debugging, and porting the 
necessary files over from my website repository 
where the model was first developed. 

\end{document}
